\documentclass[11pt,a4paper]{article}
\usepackage[margin=2.5cm]{geometry}
\usepackage{amsmath,amssymb,amsthm}
\usepackage[hidelinks]{hyperref}
\usepackage{booktabs}
\usepackage{microtype}

\newtheorem{theorem}{Theorem}
\newtheorem{corollary}{Corollary}
\newtheorem{proposition}{Proposition}
\theoremstyle{definition}
\newtheorem{definition}{Definition}
\theoremstyle{remark}
\newtheorem{remark}{Remark}
\newtheorem{conjecture}{Conjecture}

\newcommand{\Z}{\mathbb{Z}}
\newcommand{\Q}{\mathbb{Q}}
\newcommand{\C}{\mathbb{C}}
\newcommand{\mfld}[1]{\texttt{#1}}
\newcommand{\Gal}{\mathrm{Gal}}
\newcommand{\SU}{\mathrm{SU}}
\newcommand{\CS}{\mathrm{CS}}

\begin{document}

\title{The Galois--Gauge Correspondence:\\
Weyl Groups of Standard Model Forces\\
from Cusp Arithmetic of Hyperbolic Flavour Manifolds}

\author{M.~L.~Gentry%
\thanks{Independent Researcher, Seattle, WA, USA;
\texttt{drmlgentry@protonmail.com};
ORCID: 0009-0006-4550-2663}}
\date{\today}
\maketitle

\begin{abstract}
The Standard Model gauge groups $U(1)$, $\mathrm{SU}(2)$, and
$\mathrm{SU}(3)$ have Weyl groups $\{1\}$, $\mathbb{Z}/2$, and $S_3$
respectively.
We show that these three Weyl groups appear as the Galois groups of
the minimal polynomials of the cusp shape coordinates of the two
compact hyperbolic 3-manifolds of the Hyperbolic Flavour Geometry (HFG)
programme: $M_{\mathrm{PMNS}}=\mfld{m003}(-2,3)$ (the Meyerhoff manifold)
and $M_{\mathrm{CKM}}=\mfld{m006}(-5,2)$.

Specifically:
the real part $\mathrm{Re}(\tau_{\mathrm{PMNS}})=-\tfrac12$ is rational
(Galois group $\{1\}$, encoding electromagnetism);
the imaginary part $\mathrm{Im}(\tau_{\mathrm{PMNS}})=\tfrac{\sqrt3}{2}$
satisfies $x^2-\tfrac34=0$ (Galois group $\mathbb{Z}/2$, encoding the
weak force);
and the real part $\mathrm{Re}(\tau_{\mathrm{CKM}})$ satisfies the
irreducible cubic $8x^3-24x^2+28x-11$ with Galois group $S_3$
(the Weyl group of $\mathrm{SU}(3)$, encoding the strong force).
The cubic discriminant is $-3776=-2^6\times59$; the splitting field
has degree $6=|S_3|$ over $\mathbb{Q}$.

All results are verified by SnapPy and Sage computations.
This Galois--gauge correspondence provides a purely arithmetic
explanation for the three distinct gauge symmetries of the
Standard Model.
\end{abstract}

%------------------------------------------------------------------
\section{Introduction}
\label{sec:intro}

The Standard Model of particle physics has gauge group
$G_{\mathrm{SM}}=\mathrm{SU}(3)\times\mathrm{SU}(2)\times U(1)$.
The three factors govern the strong, weak, and electromagnetic forces
respectively, but no geometric or arithmetic explanation for this
specific product has been established.

In the Hyperbolic Flavour Geometry (HFG) programme~\cite{GentryUnified},
two compact hyperbolic 3-manifolds reproduce the Standard Model flavour
parameters: $M_{\mathrm{PMNS}}=\mfld{m003}(-2,3)$ (the Meyerhoff
manifold, encoding lepton mixing~\cite{GentryPMNS}) and
$M_{\mathrm{CKM}}=\mfld{m006}(-5,2)$ (encoding quark mixing).
Their cusped parents $\mfld{m003}$ and $\mfld{m006}$ carry cusp
shapes $\tau_{\mathrm{PMNS}}$ and $\tau_{\mathrm{CKM}}$ in the
upper half-plane $\mathbb{H}$.

This note identifies a precise correspondence between the
\emph{Galois groups} of the minimal polynomials of the cusp shape
coordinates and the \emph{Weyl groups} of the SM gauge factors.
This is the \textbf{Galois--gauge correspondence}.

%------------------------------------------------------------------
\section{Cusp shapes of the HFG manifolds}
\label{sec:cusps}

We use SnapPy~\cite{SnapPy} at 150-bit precision throughout.

\subsection{The lepton manifold \mfld{m003}}

The cusped manifold \mfld{m003} (SnapPy census \texttt{otet02\_00000},
two regular ideal tetrahedra with shape parameters $z_1=z_2=e^{i\pi/3}$)
has cusp shape
\begin{equation}
\tau_{\mathrm{PMNS}} = e^{i\pi/3}
= -\tfrac12 + \tfrac{\sqrt3}{2}\,i.
\label{eq:tau_pmns}
\end{equation}
This is the Eisenstein CM point with $j(\tau_{\mathrm{PMNS}})=0$.
The coordinate fields are:
\begin{align}
\mathrm{Re}(\tau_{\mathrm{PMNS}}) &= -\tfrac12 \in\mathbb{Q},
\label{eq:re_pmns}\\
\mathrm{Im}(\tau_{\mathrm{PMNS}}) &= \tfrac{\sqrt3}{2},
\quad\text{satisfying}\quad x^2-\tfrac34=0.
\label{eq:im_pmns}
\end{align}

\subsection{The quark manifold \mfld{m006}}

The cusped manifold \mfld{m006} has cusp shape (SnapPy):
\begin{equation}
\tau_{\mathrm{CKM}}
= 0.22669882\ldots + 1.46771150\ldots\, i.
\label{eq:tau_ckm}
\end{equation}
It has three tetrahedra with shapes
$z_1\approx 0.773+1.468i$ (distinct)
and $z_2=z_3\approx 0.335+0.401i$ (paired).
The relation $\mathrm{Re}(\tau_{\mathrm{CKM}})+\mathrm{Re}(z_1)=1$
holds exactly, establishing a \emph{cusp--tetrahedron complementarity}.

The Sage command \texttt{algdep(Re($z_1$), 10)} returns
\begin{equation}
p(x) = 8x^3-24x^2+28x-11,
\label{eq:cubic}
\end{equation}
verified by $p(\mathrm{Re}(z_1))<10^{-14}$.
Since $\mathrm{Re}(\tau_{\mathrm{CKM}})=1-\mathrm{Re}(z_1)$,
the minimal polynomial of $\mathrm{Re}(\tau_{\mathrm{CKM}})$ is
the mirror cubic
\begin{equation}
q(x) = p(1-x) = -8x^3-4x+1,
\label{eq:mirror}
\end{equation}
equivalently $x^3+\tfrac12 x-\tfrac18=0$ (monic depressed cubic).

%------------------------------------------------------------------
\section{The Galois--gauge correspondence}
\label{sec:galois}

\begin{theorem}[Galois--gauge correspondence]
\label{thm:main}
Let $\tau_{\mathrm{PMNS}}$ and $\tau_{\mathrm{CKM}}$ be the cusp
shapes~\eqref{eq:tau_pmns}--\eqref{eq:tau_ckm}.
Then:
\begin{enumerate}
\item $\mathrm{Re}(\tau_{\mathrm{PMNS}})=-\tfrac12$ is rational.
  Its minimal polynomial has Galois group $\{1\}$,
  the Weyl group of $U(1)$.
\item $\mathrm{Im}(\tau_{\mathrm{PMNS}})=\tfrac{\sqrt3}{2}$ satisfies
  $x^2-\tfrac34=0$.
  Its Galois group is $\mathbb{Z}/2$,
  the Weyl group of $\mathrm{SU}(2)$.
\item $\mathrm{Re}(\tau_{\mathrm{CKM}})$ satisfies the irreducible
  cubic $q(x)=x^3+\tfrac12 x-\tfrac18$ over $\mathbb{Q}$.
  Its Galois group is $S_3$,
  the Weyl group of $\mathrm{SU}(3)$.
\end{enumerate}
\end{theorem}

\begin{proof}
Items (1) and (2) are immediate from~\eqref{eq:re_pmns}
and~\eqref{eq:im_pmns}; the Galois group of $x^2-\tfrac34$ is
$\mathbb{Z}/2$ since $\tfrac34$ is not a perfect square in $\mathbb{Q}$.

For item (3): Sage computes \texttt{p.factor()} = \texttt{(8)*(irreducible
  cubic)} and \texttt{p.galois\_group()} = \emph{Transitive group number
  2 of degree~3}, which is $S_3$ (the unique non-cyclic transitive group
of degree~3).
The discriminant $\Delta=-3776<0$ confirms $S_3$ rather than $\mathbb{Z}/3$:
a cubic has Galois group $\mathbb{Z}/3$ if and only if its discriminant
is a perfect square in $\mathbb{Q}$, and $-3776$ is not.
The splitting field has degree $[K:\mathbb{Q}]=6=|S_3|$, verified by
\texttt{p.splitting\_field('a').degree()}.
\end{proof}

The Weyl group identifications used above are standard:
$W(U(1))=\{1\}$, $W(\mathrm{SU}(2))=\mathbb{Z}/2$ (the Weyl group of
the $A_1$ root system), and $W(\mathrm{SU}(3))=S_3$ (the Weyl group of
the $A_2$ root system, permuting the three fundamental weights and the
three colour charges).

\begin{corollary}
The three SM gauge factors $U(1)$, $\mathrm{SU}(2)$, $\mathrm{SU}(3)$
correspond to the three coordinate fields of the HFG cusp shapes:
the rational field $\mathbb{Q}$, the quadratic field $\mathbb{Q}(\sqrt{3})$,
and the cubic splitting field of $q(x)$ with $[\,:\mathbb{Q}]=6$.
\end{corollary}

%------------------------------------------------------------------
\section{Supporting structure}
\label{sec:support}

\subsection{The 1+2 tetrahedra split and cubic discriminant}

The discriminant $\Delta(p)=-3776=-2^6\times59$ is negative.
By the discriminant criterion, $p$ has one real root
($\mathrm{Re}(z_1)\approx0.773$) and one pair of complex conjugate roots.
This $1+2$ root structure mirrors exactly the tetrahedra structure of
\mfld{m006}: one distinct tetrahedron ($z_1$) and two identical ones
($z_2=z_3$).

\subsection{Arithmetic of the depressed cubic}

The monic depressed form $x^3+\tfrac12 x-\tfrac18=0$ has
Cardano discriminant
\begin{equation}
\Delta_C = \left(\frac{-1/8}{2}\right)^2+\left(\frac{1/2}{3}\right)^3
= \frac{1}{256}+\frac{1}{216}
= \frac{35}{1728},
\end{equation}
where $1728=12^3$ is the denominator of the $j$-function normalisation.
The numerator $35=5\times7$ and the prime $59$ in $\Delta(p)=-2^6\times59$
are new arithmetic invariants of the strong sector.

\subsection{Relation to the lepton trace field}

The cusp of \mfld{m003} has invariant trace field $K=\mathbb{Q}(\sqrt{-3})
=\mathbb{Q}(\omega)$ with $\omega=e^{i\pi/3}$ and discriminant $-3$.
The minimal polynomial of the quark cusp coordinate has discriminant
$-3776=-2^6\times59$.
The Galois group of the lepton cusp is $\mathrm{Gal}(\mathbb{Q}(\sqrt3)
/\mathbb{Q})=\mathbb{Z}/2$; of the quark cusp is $S_3$.
These two groups are related by $S_3/A_3\cong\mathbb{Z}/2$: the quark
Galois group is an extension of the lepton one by $A_3\cong\mathbb{Z}/3$.

%------------------------------------------------------------------
\section{Physical interpretation and predictions}
\label{sec:physics}

The Galois--gauge correspondence suggests a general principle:
\emph{the arithmetic complexity of the cusp shape of a hyperbolic
flavour manifold is determined by the gauge symmetry of the
corresponding force sector}.
Table~\ref{tab:main} summarises the complete picture.

\begin{table}[htbp]
\centering
\caption{The Galois--gauge correspondence for the Standard Model.
All Galois group computations are verified by Sage; cusp shapes by SnapPy.}
\label{tab:main}
\begin{tabular}{llllll}
\toprule
Force & Gauge & Weyl group & Cusp coordinate & Degree & Galois group \\
\midrule
EM    & $U(1)$          & $\{1\}$       &
  $\mathrm{Re}(\tau_{\mathrm{PMNS}})=-\tfrac12$ & 1 & $\{1\}$ \\
Weak  & $\mathrm{SU}(2)$ & $\mathbb{Z}/2$ &
  $\mathrm{Im}(\tau_{\mathrm{PMNS}})=\tfrac{\sqrt3}{2}$ & 2 & $\mathbb{Z}/2$ \\
Strong & $\mathrm{SU}(3)$ & $S_3$        &
  $\mathrm{Re}(\tau_{\mathrm{CKM}})$, cubic & 3 & $S_3$ \\
\bottomrule
\end{tabular}
\end{table}

\paragraph{Prediction: the gravitational sector.}
If gravity is encoded in a third compact hyperbolic manifold, the
Galois group of its cusp coordinate should equal the Weyl group of
the gravitational gauge group.
For Einsteinian gravity with local Lorentz symmetry $\mathrm{SO}(3,1)$,
the Weyl group is $S_4$ (the hyperoctahedral group of the $B_2/D_4$
root system in the appropriate sense); for supergravity further
structure is expected.

\paragraph{The quartic and beyond.}
The sequence $\{1\}\subset\mathbb{Z}/2\subset S_3\subset\cdots$
is the chain of Weyl groups for $U(1)\subset\mathrm{SU}(2)
\subset\mathrm{SU}(3)\subset\cdots$.
If this chain extends, the next term is $S_4$ (Weyl group of $\mathrm{SU}(4)$),
corresponding to a degree-4 cusp polynomial for a hypothetical fourth sector.
The absence of a fourth compact HFG manifold with the correct properties
would predict no fourth colour and no $\mathrm{SU}(4)$ extension.

%------------------------------------------------------------------
\section{Conclusions}
\label{sec:concl}

The Galois groups of the minimal polynomials of the HFG cusp shape
coordinates are exactly the Weyl groups of the SM gauge factors.
This Galois--gauge correspondence is:
\begin{enumerate}
\item \emph{Exact}: all computations are performed at high precision
  and verified independently by SnapPy and Sage.
\item \emph{Complete}: all three SM forces are accounted for by
  two cusp shapes.
\item \emph{Predictive}: gravity should correspond to a cusp polynomial
  with Galois group $S_4$ or larger; no fourth colour is predicted.
\item \emph{Structural}: the $1+2$ tetrahedra split of \mfld{m006},
  the negative cubic discriminant, and the $S_3$ Galois group are
  all expressions of the same underlying $N_c=3$ colour geometry.
\end{enumerate}

\section*{Data availability}
All SnapPy and Sage scripts are at
\url{https://github.com/drmlgentry/hyperbolic-flavor-scan}.

\section*{Acknowledgements}
During preparation of this manuscript the author used
Claude (claude-sonnet-4-6, Anthropic) as a research assistant.

\begin{thebibliography}{99}

\bibitem{GentryUnified}
M.~L.~Gentry,
``Hyperbolic Flavour Geometry,''
SSRN \href{https://doi.org/10.2139/ssrn.6775158}{6775158} (2026);
PTEP Paper No.\ 2605-035.

\bibitem{GentryPMNS}
M.~L.~Gentry,
``Lepton Mixing from Borel Structure of Hyperbolic Holonomy,''
SSRN \href{https://doi.org/10.2139/ssrn.6583553}{6583553} (2026);
Phys.\ Lett.\ B PLB-D-26-01448.

\bibitem{GentryWRT}
M.~L.~Gentry,
``The Gauss Polynomial and Homological Blocks of the Meyerhoff Manifold,''
preprint (2026).

\bibitem{GentryBPS}
M.~L.~Gentry,
``Lepton Masses as BPS States of a Class~S Theory on $X_0(11)$
at the Eisenstein Orbifold Point,''
preprint (2026).

\bibitem{SnapPy}
M.~Culler, N.~M.~Dunfield, M.~Goerner, J.~R.~Weeks,
SnapPy, \url{http://snappy.computop.org}.

\bibitem{GMM}
D.~Gabai, R.~Meyerhoff, P.~Milley,
J.\ Amer.\ Math.\ Soc.\ \textbf{22}, 1157 (2009).

\bibitem{Gaiotto09}
D.~Gaiotto,
``$\mathcal{N}=2$ dualities,''
JHEP \textbf{08}, 034 (2012) [arXiv:0904.2715].

\end{thebibliography}

\end{document}
